\scanpage{44}
% archive sheet number: 43
\uline{Foncteurs adjoints}

\uline{Definition 1} Soient $X,Y\in(\mathrm{Cat})$, $f:X\to Y$ et $g:Y\unsure{\text{---}|\to}X$, on dit que $(f,g)$ est un couple de foncteurs adjoints, ou que $g$ est adjoint de $f$, ou que $f$ est coadjoint de $g$, s'il existe un

\figure{\begin{tikzpicture}[x=.012cm,y=.012cm]
\draw (0,0)--(230,-2)--(48,-20)--(106,-57)--(99,13)--(104,-107)--(145,15)--(146,-84);
\draw (264,1)--(825,-23);
\draw (252,-15)--(-188,-610)--(-670,-582);
\draw (-188,-610)--(-666,-628);
\end{tikzpicture}}

\[
\begin{aligned}
&\mathscr H^{\unsure{1|{}}}(X,\underline{\operatorname{Hom}}(\gamma,G))\simeq\\
&\mathscr H^1(X,\underline{\operatorname{Ext}}^1(\gamma,G))\\
&\operatorname{Ext}^1(\gamma,
\end{aligned}
\]
\[
\begin{aligned}
&X\unsure{\text{---}|\to}\underline{\operatorname{Hom}}_{S\text{-gr}}(\gamma,G)\\
&X\times\gamma\longrightarrow G
\end{aligned}
\]
\diagram{\begin{tikzcd}[ampersand replacement=\&]
\deleted{\unsure{\gamma|g}}\,\gamma \arrow[r] \arrow[dr,dash] \& \underline{\operatorname{Hom}}_{S}(X,G)\\
S \& \overset{\unsure{P_2|P_r}}{\unsure{C|G}}
\end{tikzcd}}
\[
\operatorname{Hom}(\gamma,
\]
\[
\begin{aligned}
&\prod_{X/S}^{0}\gamma_X\\
&\quad\Vert\\
&\underline{\operatorname{Hom}}_S(X,\gamma)\unsure{-|=|\doteq}\gamma
\end{aligned}
\]
\[
e\to\gamma\to\underline{\operatorname{Hom}}_S((X,\gamma)\longrightarrow
\operatorname{Hom}(X,\unsure{a|0}\mathbin{;}\gamma)\to0
\]
